\chapter{模型扩展与工程建议}

\section{多楼层扩展}
若停车场为多层结构，可将状态扩展为 $\bm p_t=[x_t,y_t,f_t]^\T$，其中 $f_t$ 为楼层离散变量。通过电梯口、坡道口地理围栏约束楼层切换，可把二维模型扩展为“2.5 维”定位。

\section{多源融合扩展}
在 BLE 基础上加入以下传感器可进一步提升性能：
\begin{itemize}
  \item 手机 IMU（步频、航向）；
  \item UWB（局部高精度锚点）；
  \item 视觉里程计（摄像头辅助）。
\end{itemize}
融合可采用因子图框架统一表达。

\section{地图匹配扩展}
将停车位线、车道边界和出入口作为硬约束，可显著减少“穿墙点”问题。建议在优化后增加地图匹配后处理，确保轨迹符合场景语义。

\section{在线部署建议}
为满足应急应用，建议采用端云协同：
\begin{enumerate}
  \item 终端执行轻量粗定位与预处理；
  \item 边缘节点执行窗口优化与鲁棒估计；
  \item 指挥端可视化显示轨迹和可信度。
\end{enumerate}

\section{工程参数建议}
\begin{table}[H]
\centering
\caption{推荐参数范围}
\begin{tabular}{lll}
\toprule
参数 & 建议范围 & 说明 \\
\midrule
窗口长度 $w$ & 6--10 & 平衡时延与平滑性 \\
Huber 阈值 $\delta$ & 2--4 m & 异常抑制强度 \\
$\lambda_s$ & 0.1--1.0 & 速度平滑权重 \\
$\lambda_h$ & 0.01--0.2 & 转向平滑权重 \\
\bottomrule
\end{tabular}
\end{table}

\section{风险与应对}
部署阶段的主要风险包括信标电量衰减、位置漂移和无线干扰。建议建立巡检机制，周期性校准信标功率并自动检测异常锚点。
